% GENERATED FILE. Edit canonical lessons, then run npm run generate:books.
% canonical-source-hash: fnv1a64:6b92a7ed800bc7d5
% canonical-lessons: HI-C493-madhumeh, HI-C493-dama, HI-C493-lakva, HI-C493-kaimsar, HI-C493-jakarna

\chapter{Long Illnesses}
\label{ch:hi-long-illnesses}
\hlchaptermodality{493}

\begin{chapteropening}
\textbf{By the end of this chapter:} \emph{I can talk about an illness that lasts.}

Complete the last lesson of chapter 493: I can talk about an illness that lasts.
\end{chapteropening}

\section[\texorpdfstring{\hi{मधुमेह} (madhumeh)}{मधुमेह (madhumeh)}]{\hi{मधुमेह} (madhumeh) — diabetes}

\label{lesson:HI-C493-madhumeh}



\begin{quote}
Before the new one: say the Hindi for ayurveda, then the Hindi for vigour.
\end{quote}

\subsection*{You'll want to know: \hi{मधुमेह}}
\textbf{\hi{मधुमेह}} — \emph{madhumeh} — \textquotedblleft{}diabetes\textquotedblright{}.

\begin{quote}
\hi{मधुमेह} \hi{के} \hi{मरीज़} \hi{को} \hi{मीठा} \hi{कम} \hi{खाना} \hi{चाहिए।}
\end{quote}

\begin{quote}
\emph{madhumeh ke marīz ko mīṭhā kam khānā cāhiye.}
\end{quote}

\begin{quote}
Someone with diabetes should eat fewer sweet things.
\end{quote}

Masculine.

\begin{cousinweb}
\textbf{\hi{मधुमेह}} (\emph{madhumeh}) is a tatsama, borrowed from Sanskrit \emph{madhumeha}, \textquotedblleft{}honey-like urine, diabetes\textquotedblright{} (Monier-Williams), built on \emph{madhu}, \textquotedblleft{}honey\textquotedblright{}, and \emph{meha}, \textquotedblleft{}urine\textquotedblright{}: old doctors knew the illness by its sweetness. Speech often says \emph{śugar}.
\end{cousinweb}

\begin{grammarlens}[title={What you've built}]
One more word for describing what you lived and felt.
\end{grammarlens}

\subsection*{Your turn}
\begin{itemize}
\raggedright
  \item \emph{Say it:} \emph{madhumeh}
  \item \emph{Say it:} \emph{madhumeh}, once more
  \item \emph{Say it:} \emph{madhumeh ke marīz ko mīṭhā kam khānā cāhiye.} — Someone with diabetes should eat fewer sweet things.
  \item \emph{Recall:} say the Hindi for ayurveda, then the Hindi for vigour, then say \emph{madhumeh} again
\end{itemize}

\begin{tcolorbox}[breakable,colback=teal!4,colframe=teal!35!black,title={Before you move on}]
What does \hi{मधुमेह} mean? (\textquotedblleft{}Diabetes\textquotedblright{}.) Say it once more.
\end{tcolorbox}
\section[\texorpdfstring{\hi{दमा} (damā)}{दमा (damā)}]{\hi{दमा} (damā) — asthma}

\label{lesson:HI-C493-dama}



\begin{quote}
Before the new one: say the Hindi for vigour, then the Hindi for diabetes.
\end{quote}

\subsection*{You'll want to know: \hi{दमा}}
\textbf{\hi{दमा}} — \emph{damā} — \textquotedblleft{}asthma\textquotedblright{}.

\begin{quote}
\hi{दादी} \hi{को} \hi{दमा} \hi{है}, \hi{इसलिए} \hi{वह} \hi{ठंडी} \hi{हवा} \hi{से} \hi{दूर} \hi{रहती} \hi{हैं।}
\end{quote}

\begin{quote}
\emph{dādī ko damā hai, isliye vah ṭhaṇḍī havā se dūr rahtī haiṁ.}
\end{quote}

\begin{quote}
Grandmother has asthma, so she keeps away from cold air.
\end{quote}

Masculine.

\begin{cousinweb}
\textbf{\hi{दमा}} (\emph{damā}) is usually traced to Persian \emph{dama}, \textquotedblleft{}asthma\textquotedblright{}, a word beside \emph{dam}, \textquotedblleft{}breath\textquotedblright{}: the illness of the breath.
\end{cousinweb}

\begin{grammarlens}[title={What you've built}]
One more word for describing what you lived and felt.
\end{grammarlens}

\subsection*{Your turn}
\begin{itemize}
\raggedright
  \item \emph{Say it:} \emph{damā}
  \item \emph{Say it:} \emph{damā}, once more
  \item \emph{Say it:} \emph{dādī ko damā hai, isliye vah ṭhaṇḍī havā se dūr rahtī haiṁ.} — Grandmother has asthma, so she keeps away from cold air.
  \item \emph{Recall:} say the Hindi for vigour, then the Hindi for diabetes, then say \emph{damā} again
\end{itemize}

\begin{tcolorbox}[breakable,colback=teal!4,colframe=teal!35!black,title={Before you move on}]
What does \hi{दमा} mean? (\textquotedblleft{}Asthma\textquotedblright{}.) Say it once more.
\end{tcolorbox}
\section[\texorpdfstring{\hi{लकवा} (lakvā)}{लकवा (lakvā)}]{\hi{लकवा} (lakvā) — paralysis; a stroke}

\label{lesson:HI-C493-lakva}



\begin{quote}
Before the new one: say the Hindi for diabetes, then the Hindi for asthma.
\end{quote}

\subsection*{You'll want to know: \hi{लकवा}}
\textbf{\hi{लकवा}} — \emph{lakvā} — \textquotedblleft{}paralysis; a stroke\textquotedblright{}.

\begin{quote}
\hi{लकवे} \hi{के} \hi{बाद} \hi{दादा} \hi{का} \hi{एक} \hi{हाथ} \hi{कमज़ोर} \hi{है।}
\end{quote}

\begin{quote}
\emph{lakve ke bād dādā kā ek hāth kamzor hai.}
\end{quote}

\begin{quote}
Since his stroke, my grandfather has one weak hand.
\end{quote}

Masculine.

\begin{cousinweb}
\textbf{\hi{लकवा}} (\emph{lakvā}) is usually traced to Arabic \emph{laqwa}, \textquotedblleft{}facial palsy\textquotedblright{}, which reached Hindi through Persian and widened to any paralysis.
\end{cousinweb}

\begin{grammarlens}[title={What you've built}]
One more word for describing what you lived and felt.
\end{grammarlens}

\subsection*{Your turn}
\begin{itemize}
\raggedright
  \item \emph{Say it:} \emph{lakvā}
  \item \emph{Say it:} \emph{lakvā}, once more
  \item \emph{Say it:} \emph{lakve ke bād dādā kā ek hāth kamzor hai.} — Since his stroke, my grandfather has one weak hand.
  \item \emph{Recall:} say the Hindi for diabetes, then the Hindi for asthma, then say \emph{lakvā} again
\end{itemize}

\begin{tcolorbox}[breakable,colback=teal!4,colframe=teal!35!black,title={Before you move on}]
What does \hi{लकवा} mean? (\textquotedblleft{}Paralysis; a stroke\textquotedblright{}.) Say it once more.
\end{tcolorbox}
\section[\texorpdfstring{\hi{कैंसर} (kaiṁsar)}{कैंसर (kaiṁsar)}]{\hi{कैंसर} (kaiṁsar) — cancer}

\label{lesson:HI-C493-kaimsar}



\begin{quote}
Before the new one: say the Hindi for asthma, then the Hindi for paralysis.
\end{quote}

\subsection*{You'll want to know: \hi{कैंसर}}
\textbf{\hi{कैंसर}} — \emph{kaiṁsar} — \textquotedblleft{}cancer\textquotedblright{}.

\begin{quote}
\hi{डाक्टर} \hi{ने} \hi{कहा} \hi{कि} \hi{दादा} \hi{को} \hi{कैंसर} \hi{नहीं} \hi{है।}
\end{quote}

\begin{quote}
\emph{ḍākṭar ne kahā ki dādā ko kaiṁsar nahīṃ hai.}
\end{quote}

\begin{quote}
The doctor said my grandfather doesn't have cancer.
\end{quote}

Masculine.

\begin{cousinweb}
\textbf{\hi{कैंसर}} (\emph{kaiṁsar}) is borrowed from English \emph{cancer}.
\end{cousinweb}

\begin{grammarlens}[title={What you've built}]
One more word for describing what you lived and felt.
\end{grammarlens}

\subsection*{Your turn}
\begin{itemize}
\raggedright
  \item \emph{Say it:} \emph{kaiṁsar}
  \item \emph{Say it:} \emph{kaiṁsar}, once more
  \item \emph{Say it:} \emph{ḍākṭar ne kahā ki dādā ko kaiṁsar nahīṃ hai.} — The doctor said my grandfather doesn't have cancer.
  \item \emph{Recall:} say the Hindi for asthma, then the Hindi for paralysis, then say \emph{kaiṁsar} again
\end{itemize}

\begin{tcolorbox}[breakable,colback=teal!4,colframe=teal!35!black,title={Before you move on}]
What does \hi{कैंसर} mean? (\textquotedblleft{}Cancer\textquotedblright{}.) Say it once more.
\end{tcolorbox}
\section[\texorpdfstring{\hi{जकड़ना} (jakaṛnā)}{जकड़ना (jakaṛnā)}]{\hi{जकड़ना} (jakaṛnā) — to stiffen; to grip tight}

\label{lesson:HI-C493-jakarna}



\begin{quote}
Before the new one: say the Hindi for paralysis, then the Hindi for cancer.
\end{quote}

\subsection*{You'll want to know: \hi{जकड़ना}}
\textbf{\hi{जकड़ना}} — \emph{jakaṛnā} — \textquotedblleft{}to stiffen; to grip tight\textquotedblright{}.

\begin{quote}
\hi{ठंड} \hi{में} \hi{दादी} \hi{के} \hi{घुटने} \hi{जकड़} \hi{जाते} \hi{हैं।}
\end{quote}

\begin{quote}
\emph{ṭhaṇḍ meṁ dādī ke ghuṭne jakaṛ jāte haiṁ.}
\end{quote}

\begin{quote}
In the cold my grandmother's knees stiffen up.
\end{quote}

Past (it …): \hi{वह} \hi{घुटना} \hi{जकड़} \hi{गया}, \emph{vah ghuṭnā jakaṛ gayā}.

\begin{cousinweb}
\textbf{\hi{जकड़ना}} (\emph{jakaṛnā}) is inherited: CDIAL lists Hindi \emph{jakaṛnā}, \textquotedblleft{}to tighten, bind\textquotedblright{}, under Sanskrit \emph{yatá}, \textquotedblleft{}held in, restrained\textquotedblright{} (CDIAL 10400). A stiff joint feels bound tight.
\end{cousinweb}

\begin{grammarlens}[title={What you've built}]
One more word for describing what you lived and felt.
\end{grammarlens}

\subsection*{Your turn}
\begin{itemize}
\raggedright
  \item \emph{Say it:} \emph{jakaṛnā}
  \item \emph{Say it:} \emph{jakaṛnā}, once more
  \item \emph{Say it:} \emph{ṭhaṇḍ meṁ dādī ke ghuṭne jakaṛ jāte haiṁ.} — In the cold my grandmother's knees stiffen up.
  \item \emph{Recall:} say the Hindi for paralysis, then the Hindi for cancer, then say \emph{jakaṛnā} again
\end{itemize}

\begin{tcolorbox}[breakable,colback=teal!4,colframe=teal!35!black,title={Before you move on}]
What does \hi{जकड़ना} mean? (\textquotedblleft{}To stiffen; to grip tight\textquotedblright{}.) Say it once more.
\end{tcolorbox}
